\documentclass[11pt,letterpaper]{article}
% Explanatory companion, not the workshop's five-page submission template.
\usepackage[margin=0.9in]{geometry}
\usepackage[T1]{fontenc}
\usepackage{lmodern}
\usepackage{amsmath,amssymb}
\usepackage{newunicodechar}
\newunicodechar{−}{\ensuremath{-}}
\usepackage{graphicx}
\usepackage{booktabs,longtable,array,calc}
\usepackage{microtype}
\usepackage{xcolor}
\definecolor{linkblue}{HTML}{145A7A}
\usepackage{hyperref}
\hypersetup{colorlinks=true,linkcolor=linkblue,urlcolor=linkblue,
  pdftitle={BrainBodyFM 2026: submission-readiness
guide},pdfauthor={},pdfsubject={BrainBodyFM tutorial and evidence review}}
\usepackage{caption}
\captionsetup{font=small,labelfont=bf}
% Pandoc 3.10 uses LTcaptype=none for uncaptioned longtables.
% The older ltcaption bundled with Tectonic still expects a counter.
\newcounter{none}
\usepackage{float}
\usepackage{needspace}
\floatplacement{figure}{H}
\usepackage{etoolbox}
\AtBeginEnvironment{longtable}{\small}
\setlength{\emergencystretch}{3em}
\setlength{\parindent}{0pt}
\setlength{\parskip}{0.55em}
\setcounter{secnumdepth}{0}
\providecommand{\tightlist}{\setlength{\itemsep}{0pt}\setlength{\parskip}{0pt}}
\providecommand{\pandocbounded}[1]{#1}
\title{\LARGE BrainBodyFM 2026: submission-readiness
guide\\[0.45em]\large Current two-study manuscript and remaining
decisions}
\author{}
\date{\small Updated September 5, 2026}
\begin{document}
\maketitle
\section{Current position}\label{current-position}

The revised short paper combines a source-held-out classification pilot
with a completed laterality study. Its central contribution is an
empirical evaluation of what learned gait features preserve, how
reflection augmentation changes token consistency, and what an
explicitly imposed sign constraint can guarantee. The laterality study
supplies the stronger controlled evidence; the classification pilot
provides supporting context.

The editable submission source is \url{bbfm2026_paper_draft.md}, with
generated \href{bbfm2026_paper_draft.tex}{LaTeX} and
\href{bbfm2026_paper_draft.pdf}{PDF}. The
\href{explanatory_paper.pdf}{explanatory paper} provides the longer
methods tutorial and evidence assessment. It is an author resource
rather than the main submission file.

Submission and release remain blocked by the laterality project's
recorded governance gate. Ethics determination, data-use review, and
derived-pose release review are all unresolved in
\href{../../neurips-laterality/governance/status.json}{governance/status.json}.
The implemented rule requires each review to be resolved with an
internal reference and date. This is a documented project condition, not
a claim about an automatic rejection rule of the workshop. A weighted
readiness score cannot override it.

\section{Two studies, with different
evidence}\label{two-studies-with-different-evidence}

{\def\LTcaptype{none} % do not increment counter
\begin{longtable}[]{@{}
  >{\raggedright\arraybackslash}p{(\linewidth - 4\tabcolsep) * \real{0.3333}}
  >{\raggedright\arraybackslash}p{(\linewidth - 4\tabcolsep) * \real{0.3333}}
  >{\raggedright\arraybackslash}p{(\linewidth - 4\tabcolsep) * \real{0.3333}}@{}}
\toprule\noalign{}
\begin{minipage}[b]{\linewidth}\raggedright
Property
\end{minipage} & \begin{minipage}[b]{\linewidth}\raggedright
Study A: classification pilot
\end{minipage} & \begin{minipage}[b]{\linewidth}\raggedright
Study B: laterality
\end{minipage} \\
\midrule\noalign{}
\endhead
\bottomrule\noalign{}
\endlastfoot
Accepted cohort & 639 sequences / 97 sources & 625 sequences / 93
sources \\
Evaluation & Outer fold 0, seed 42 & Five outer folds, five seeds, two
training variants \\
Source counts & 59 train, 18 validation, 20 test & 74--75 outer train,
18--19 outer test per fold \\
Main comparison & Learned features, raw pose statistics, missingness &
Trained versus recorded initial features, reflection augmentation,
odd/even readouts \\
Current artifacts & Saved notebook outputs; current
checkpoint/evaluation bundles absent & Cohort, splits, 50 trained
checkpoints, and 100,000 prediction rows across 16 lanes inspected \\
\end{longtable}
}

The cohorts overlap and use different eligibility rules and models.
Their counts should not be added, and Study B does not replicate Study
A's classifier experiment. In both studies, source-video separation is
weaker than person separation because one individual may appear in
multiple uploads.

Study A's dated census begins with 666 annotated sequences from 103
uploads. Its September 4 metadata gate retains 657 sequences/100
sources, the media-span gate 655/98, and pose quality control 639/97.
Source roles are frozen before the later exclusions. These figures
describe the saved acquisition snapshot, not current video availability.

Study B starts with a frozen inventory of 642 pose archives and retains
625/93 after quality and target-computability checks. The local
verification loaded its recorded lineage and predictions. Independent
in-memory recomputation of the primary predictive and token-equivariance
bootstrap tables reproduced the saved results within \(10^{-16}\). That
checks the stored evidence and report calculations; no independent
retraining was performed.

\section{The results to emphasize}\label{the-results-to-emphasize}

Study A's source-level macro-F1 is 0.440513 for raw pose statistics,
0.292424 for learned features, and 0.251111 for missingness. All three
lanes miss all three stroke-annotated test sources. The comparison
concerns 20 test sources in one fold and has no matched untrained
encoder, so it does not isolate a general effect of pretraining.

Study B provides the following paired contrasts. Intervals are pointwise
95\% source-bootstrap intervals, conditional on the fitted models and
fixed split.

{\def\LTcaptype{none} % do not increment counter
\begin{longtable}[]{@{}
  >{\raggedright\arraybackslash}p{(\linewidth - 4\tabcolsep) * \real{0.3333}}
  >{\raggedright\arraybackslash}p{(\linewidth - 4\tabcolsep) * \real{0.3333}}
  >{\raggedright\arraybackslash}p{(\linewidth - 4\tabcolsep) * \real{0.3333}}@{}}
\toprule\noalign{}
\begin{minipage}[b]{\linewidth}\raggedright
Contrast
\end{minipage} & \begin{minipage}[b]{\linewidth}\raggedright
Estimate {[}95\% interval{]}
\end{minipage} & \begin{minipage}[b]{\linewidth}\raggedright
Supported reading
\end{minipage} \\
\midrule\noalign{}
\endhead
\bottomrule\noalign{}
\endlastfoot
Vanilla trained minus initial token error & +0.03055 {[}0.01576,
0.04763{]} & Training increases error under the specified token
action \\
Reflection-augmented minus vanilla token error & −0.00843 {[}−0.01020,
−0.00687{]} & Augmentation improves this geometric consistency
measure \\
Native learned minus initial predictive \(R^2\) & −0.01798 {[}−0.03851,
0.00248{]} & A predictive benefit is unestablished; equivalence is also
unestablished \\
Constructed odd learned minus initial \(R^2\) & −0.05874 {[}−0.09549,
−0.01740{]} & Trained features perform worse through this constrained
readout \\
Reflection-augmented minus vanilla native \(R^2\) & +0.00408
{[}−0.00556, 0.01277{]} & Predictive improvement remains uncertain \\
\end{longtable}
}

An odd feature construction and origin-preserving linear readout enforce
output sign reversal for any encoder, including its initialization. The
saved study verifies that property for every seed. Its value is as a
controlled transformation rule; useful learned information still
requires a predictive comparison. The novelty is the specific empirical
separation of these outcomes, with matched controls, rather than new
symmetry algebra or a discovery of anatomical laterality.

Study B's primary predictive score uses sequence rows weighted by
inverse source size. For each seed, it pools out-of-fold predictions
before calculating \(R^2\), then averages the five seed scores. It does
not average source targets first or score a prediction ensemble. Study
A's temporal diagnostic does average targets and predictions within
source before computing its separate \(R^2\).

\section{Method descriptions that must remain
accurate}\label{method-descriptions-that-must-remain-accurate}

Study A uses four-frame coordinate averages, a pooled-context MLP
predictor, and SmoothL1 latent prediction plus variance and covariance
penalties weighted 0.10 and 0.01. Its optional 0.10 condition
cross-entropy term is disabled. Labels still determine the cumulative
normal-first curriculum. Classifier selection uses validation sources,
followed by a scaler/classifier refit on all 77 development sources. The
encoder receives gradient updates from the 59 training sources only.

Study B flattens four-frame coordinate patches into 96-dimensional
tokens, uses a Transformer predictor, and trains with centered latent
cross-entropy plus a 0.05-weighted VICReg-style objective. Five folds,
five seeds, and two variants give 50 trained fits. These implementations
should not share one undifferentiated loss equation or configuration
description.

Study A's temporal and drift notebooks interpolate short gaps, unlike
its training preparation. Its classifier also averages probabilities at
test time while validation uses source-mean features. Those limitations
remain even after their descriptions are corrected. Normal-anchor cosine
measures same-clip geometric change; functional forgetting and a
successful consolidation repair remain untested. The forecasting branch
lacks a separately trained checkpoint and requires a prefix-only context
and baseline implementation before a causal forecasting claim is
possible.

\section{Internal readiness
assessment}\label{internal-readiness-assessment}

Scores run from 1 to 5, with 3 representing an adequate limited
contribution and 4 a strong one. The revised index assesses the
integrated paper with the factual corrections applied and the evidence
limitations retained. It is an author-side heuristic, not a workshop
rubric or acceptance probability.

{\def\LTcaptype{none} % do not increment counter
\begin{longtable}[]{@{}lrr@{}}
\toprule\noalign{}
Dimension & Weight & Revised score / 5 \\
\midrule\noalign{}
\endhead
\bottomrule\noalign{}
\endlastfoot
Workshop fit & 15\% & 4.0 \\
Contribution and novelty & 15\% & 3.0 \\
Methods correctness & 20\% & 4.0 \\
Empirical strength & 20\% & 3.5 \\
Reproducibility and traceability & 15\% & 3.5 \\
Claim discipline & 5\% & 4.5 \\
Clarity and presentation & 5\% & 4.0 \\
Data-use reporting and submission readiness & 5\% & 2.0 \\
\end{longtable}
}

Movement-representation evaluation gives the paper a direct workshop
connection, while its methods remain established. The repeated
laterality study and matched controls support the methods score;
empirical strength remains limited by one observational dataset and no
external person-level test. Reproducibility is scored for the whole
paper, including the missing pilot bundles. The claim and clarity scores
assume the revised two-study framing is retained. Unresolved governance
keeps the final dimension low.

With percentage weights \(w_j\), the index is
\(\sum_j w_j s_j/5=72/100\). The historical uncorrected pilot-only draft
scored 48/100; a hypothetical corrected pilot without additional
evidence scored approximately 57/100. The increase to 72 reflects the
integrated existing laterality evidence and verification. It does not
mean the pilot gained additional folds or that acceptance odds are 72\%.

The scientific case is now more credible as a bounded methodological
workshop contribution. Its novelty remains incremental, and no neural
measurements, broad foundation-model transfer, clinical validation, or
interactive control results are available. Independently of this
assessment, the governance gate remains blocked.

\section{Remaining priorities}\label{remaining-priorities}

\begin{enumerate}
\def\labelenumi{\arabic{enumi}.}
\tightlist
\item
  Resolve the three recorded governance reviews through the responsible
  authors or institution and record the required references and dates.
  Do not mark them resolved on the basis of manuscript edits.
\item
  Preserve the verified short-paper format in any final revision. The
  September 5 build has five main-text pages, two appendix pages, and
  one reference page in the shared NeurIPS 2026 double-blind workshop
  style. All eight pages were visually inspected; author identity is
  absent from the PDF metadata, and the vector figure uses labels of at
  least 9.3 pt at the official 5.5-inch text width. Recheck the PDF
  after any further edit.
\item
  Check that every abstract, table, and caption distinguishes the two
  cohorts and estimands. Preserve the token-error benefit of
  augmentation, the inconclusive native predictive contrast, and the
  negative constructed learned-minus-initial result.
\item
  Confirm that any linked reviewer-facing artifact extract complies with
  the recorded release determination. The local evidence contains
  source-linked records and is not automatically authorized for public
  redistribution.
\item
  Recover Study A's original bundles if available. If they remain
  unavailable, retain the explicit saved-output limitation and avoid
  upgrading the pilot's evidence status.
\item
  Verify the live submission portal and required author declarations
  before upload. The published deadline is September 5, 2026 AoE; the
  \href{https://brainbodyfm-workshop.github.io/call-for-papers.html}{call
  for papers} and
  \href{https://openreview.net/group?id=NeurIPS.cc/2026/Workshop/BrainBodyFM}{OpenReview
  venue} are the relevant submission references.
\end{enumerate}

New architecture changes, pilot preprocessing fixes, or forecasting
experiments would create new evidence. They should not inherit the
current scores without rerunning the affected pipeline. The completed
laterality results can support the present paper while these broader
experiments remain follow-up work.
\end{document}
